% !TeX root = ../../../manual.tex
\chapter{Tables}
\label{ch:tables}

\section{booktabs}

The \texttt{booktabs} package provides professional-quality horizontal rules.
OmniLaTeX loads it by default. The booktabs philosophy discourages vertical
lines and encourages generous spacing.

\begin{verbatim}
\begin{tabular}{lcc}
    \toprule
    Parameter   & Value  & Unit \\
    \midrule
    Temperature & 298.15 & K    \\
    Pressure    & 101.3  & kPa  \\
    \addlinespace
    Density     & 1.225  & kg/m$^{3}$ \\
    \bottomrule
\end{tabular}
\end{verbatim}

\begin{longtblr}[
    caption = {booktabs rule commands},
    label = {tab:booktabs-rules},
]{colspec = {l l}, rowhead = 1}
    \toprule
    Command          & Purpose \\
    \midrule
    \texttt{\textbackslash toprule}    & Heavy rule at the top of the table \\
    \texttt{\textbackslash midrule}    & Medium rule below the header \\
    \texttt{\textbackslash bottomrule} & Heavy rule at the bottom of the table \\
    \texttt{\textbackslash addlinespace} & Insert extra vertical space between rows \\
    \bottomrule
\end{longtblr}

Avoid \verb|\hline| and vertical lines (\verb|{|}) in booktabs tables.

\section{siunitx S-Columns}

The \texttt{siunitx} package provides the \texttt{S} column type for
automatic alignment of numbers by their decimal point:

\begin{verbatim}
\begin{tabular}{l S[table-format=3.2] S[table-format=3.2]}
    \toprule
    {Material} & {Length (mm)} & {Mass (g)} \\
    \midrule
    Steel      & 125.40        & 492.10     \\
    Aluminum   &  98.70        &  33.55     \\
    Copper     & 100.00        & 893.30     \\
    \bottomrule
\end{tabular}
\end{verbatim}

Enclose column headers in braces (\texttt{\{...\}}) to prevent siunitx from
parsing them as numeric values. The \texttt{table-format} key specifies the
expected number of digits before and after the decimal separator.

\section{tabularray}

OmniLaTeX recommends the \texttt{tabularray} package for complex tables. It
provides a modern key-value interface, long table support, and built-in
row coloring.

\begin{verbatim}
\begin{longtblr}[
    caption = {Experimental results},
    label   = {tab:results},
    remark  = {Note},
]{colspec = {l S[table-format=2.2] S[table-format=2.2]},
    rowhead = 1,
    row{odd} = {gray9},
}
    \toprule
    Sample & {Voltage (V)} & {Current (A)} \\
    \midrule
    A      & 3.14          & 0.42          \\
    B      & 5.00          & 1.07          \\
    C      & 2.71          & 0.88          \\
    \bottomrule
\end{longtblr}
\end{verbatim}

Variable-width columns use \texttt{l}, \texttt{c}, \texttt{r}, or explicit
widths: \texttt{X[2,c]} for a centered column twice the default width.

\section{Multi-Column Cells}

Span cells across multiple columns with \verb|\multicolumn|:

\begin{verbatim}
\begin{tabular}{lccc}
    \toprule
    Experiment & \multicolumn{2}{c}{Results} & Status \\
    \cmidrule(lr){2-3}
    Run 1      & 0.95 & 0.97 & Pass \\
    Run 2      & 0.82 & 0.79 & Fail \\
    \bottomrule
\end{tabular}
\end{verbatim}

Use \verb|\cmidrule(lr){2-3}| for partial horizontal rules that span only the
merged columns. The \texttt{(lr)} trims shorten the rule at both ends for a
cleaner appearance.

\section{Rotated Cells}

Rotate text within table cells using \verb|\multirotatecell|:

\begin{verbatim}
\multirotatecell{3}{Rotated text spanning three rows}
\end{verbatim}

This is particularly useful for column headers with long labels in wide
tables. The first argument specifies the number of rows the cell should span.

\section{Table Footnotes}

Add footnotes to tabularray tables with \verb|\TblrNote| and
\verb|\TblrRemark|:

\begin{verbatim}
\begin{longtblr}[remark = {Notes}]{colspec = {l S}, remark = {Note}}
    \toprule
    Entry & {Value} \\
    \midrule
    Alpha & 1.23\TblrNote{a} \\
    Beta  & 4.56\TblrNote{b} \\
    Gamma & 7.89 \\
    \bottomrule
\end{longtblr}
\TblrNote{a}{Measured at 25\,°C.}
\TblrNote{b}{Average of three trials.}
\TblrRemark{c}{Data collected in 2025.}
\end{verbatim}

Notes are collected below the table and referenced by key.

\section{Table Styles}

OmniLaTeX defines several tabularray templates for consistent table styling:

\begin{longtblr}[
    caption = {OmniLaTeX tabularray templates},
    label = {tab:tblr-templates},
]{colspec = {l l}, rowhead = 1}
    \toprule
    Template      & Purpose \\
    \midrule
    contfoot-text & Table with continued footer text across pages \\
    conthead-text & Table with continued header text across pages \\
    caption-sep   & Enforced vertical separation between caption and table body \\
    \bottomrule
\end{longtblr}

Apply a template with the \texttt{theme} key:

\begin{verbatim}
\begin{longtblr}[theme = caption-sep, caption = {...}]
    ...
\end{longtblr}
\end{verbatim}

These templates are defined in the OmniLaTeX core style files and can be
extended or overridden in your own preamble.
